
\documentclass{article}
\usepackage{colortbl}
\usepackage{makecell}
\usepackage{multirow}
\usepackage{supertabular}

\begin{document}

\newcounter{utterance}

\centering \large Interaction Transcript for game `eqbench', experiment `eq\_bench', episode 128 with reflexion{-}Qwen3.5{-}27B.
\vspace{24pt}

{ \footnotesize  \setcounter{utterance}{1}
\setlength{\tabcolsep}{0pt}
\begin{supertabular}{c@{$\;$}|p{.15\linewidth}@{}p{.15\linewidth}p{.15\linewidth}p{.15\linewidth}p{.15\linewidth}p{.15\linewidth}}
    \# & \multicolumn{2}{c}{Player} && \multicolumn{2}{c}{Game Master} \\
    \hline

    \theutterance \stepcounter{utterance}  
    & & & \multicolumn{4}{p{0.6\linewidth}}{
        \cellcolor[rgb]{0.9,0.9,0.9}{
            \makecell[{{p{\linewidth}}}]{
                \texttt{\tiny{[P1$\langle$GM]}}
                \texttt{You have attempted this same task before and did not succeed. These are the plans you wrote after those attempts.} \\
\texttt{Trial \#1: New plan: My previous attempt failed because I assigned a high score to "Annoyance" (8) and a non{-}zero score to "Resentment" (4), likely misinterpreting Olivia's defensive and dismissive tone as active anger rather than guarded skepticism. Olivia's dialogue ("I suppose," "not here for a philosophical discussion," "makes things more complicated") indicates she is resistant to the doctor's probing but not necessarily hostile or deeply resentful. The most accurate emotional state is likely a mix of "Indifference" (disinterest in the topic) and mild "Annoyance" at the persistence, but "Resentment" and "Acceptance" should be near zero. Next time, I will re{-}evaluate the dialogue to distinguish between active hostility and defensive disengagement. I will assign a high score to "Indifference" (e.g., 7) to reflect her dismissal of the topic, a moderate score to "Annoyance" (e.g., 4) for her irritation at the doctor's persistence, and near{-}zero scores to "Acceptance" (0) and "Resentment" (1), ensuring the scores accurately reflect her guarded, non{-}committal stance rather than assuming deep{-}seated anger.} \\
\\ 
\texttt{Trial \#2: New plan: My previous attempt failed because I overestimated "Indifference" (7) and underestimated "Annoyance" (6), misinterpreting Olivia's sharp, defensive retorts as mere disinterest. Her lines ("I'm not here for a philosophical discussion," "Or it just makes things more complicated") are active pushbacks against the doctor's persistence, indicating clear irritation rather than a lack of caring. "Indifference" implies she doesn't care about the interaction, but she is actively engaging to shut it down. Next time, I will prioritize "Annoyance" as the dominant emotion (e.g., 8) to reflect her frustration with the doctor's circular logic. I will lower "Indifference" significantly (e.g., 2) to acknowledge she is engaged but negatively. "Resentment" should remain low (1) as there is no evidence of deep{-}seated bitterness, and "Acceptance" should be 0. The key is to distinguish between "I don't care" (Indifference) and "I am irritated by this" (Annoyance).} \\
\\ 
\texttt{Trial \#3: New plan: My previous attempts failed because I oscillated between overestimating "Indifference" and overestimating "Annoyance" without a consistent analytical framework. In Trial 1, I misread her active pushback as disinterest; in Trial 2, I correctly identified the irritation but likely failed to account for the nuance of her specific phrasing which suggests a defensive barrier rather than pure anger. The dialogue shows Olivia actively rejecting the doctor's premise ("not here for a philosophical discussion") and countering his logic ("makes things more complicated"). This is not indifference (she is engaging to stop him) nor deep resentment (no history of betrayal is cited, just current friction). The most accurate interpretation is that she is **Annoyed** by the circular, unhelpful nature of the conversation. My new strategy is to strictly define the emotions: "Indifference" is 0 because she is actively arguing; "Acceptance" is 0 because she rejects the premise; "Resentment" is low (1{-}2) as it's situational, not deep{-}seated; and "Annoyance" is the dominant score (8{-}9) reflecting her frustration with the doctor's persistence. I will output scores that reflect this hierarchy: Annoyance high, Resentment low, Indifference and Acceptance zero.} \\
\\ 
\\ 
\\ 
\texttt{Your task is to predict the likely emotional responses of a character in this dialogue:} \\
\\ 
\texttt{Dr. Fleischer: You've been doing well, Olivia. Your writing's been improving, hasn't it?} \\
\texttt{Olivia: I suppose. But that's not why I'm here, is it?} \\
\texttt{Dr. Fleischer: No, it's not. But it's all connected, isn't it? Your writing, your past, your anxieties...} \\
\texttt{Olivia: I'm not here for a philosophical discussion, Dr. Fleischer.} \\
\texttt{Dr. Fleischer: Of course, you're not. But sometimes, understanding the connections helps us untangle the knots, wouldn't you agree?} \\
\texttt{Olivia: Or it just makes things more complicated.} \\
\texttt{Dr. Fleischer: Perhaps. But isn't it better to face the complexity than to ignore it?} \\
\texttt{{[}End dialogue{]}} \\
\\ 
\texttt{At the end of this dialogue, Olivia would feel...} \\
\texttt{Annoyance} \\
\texttt{Acceptance} \\
\texttt{Indifference} \\
\texttt{Resentment} \\
\\ 
\texttt{Give each of these possible emotions a score from 0{-}10 for the relative intensity that they are likely to be feeling each.} \\
\\ 
\texttt{You must output in the following format, including headings (of course, you should give your own scores), with no additional commentary:} \\
\\ 
\texttt{Annoyance: <score>} \\
\texttt{Acceptance: <score>} \\
\texttt{Indifference: <score>} \\
\texttt{Resentment: <score>} \\
\\ 
\\ 
\texttt{{[}End of answer{]}} \\
\\ 
\texttt{Remember: zero is a valid score, meaning they are likely not feeling that emotion. You must score at least one emotion > 0.} \\
\\ 
\texttt{Your answer:} \\
            }
        }
    }
     \\ \\

    \theutterance \stepcounter{utterance}  
    & \multicolumn{4}{p{0.6\linewidth}}{
        \cellcolor[rgb]{0.561,0.737,0.561}{
            \makecell[{{p{\linewidth}}}]{
                \texttt{\tiny{[P1$\rangle$GM]}}
                \texttt{Annoyance: 8} \\
\texttt{Acceptance: 0} \\
\texttt{Indifference: 1} \\
\texttt{Resentment: 2} \\
            }
        }
    }
    & & \\ \\

    \theutterance \stepcounter{utterance}  
    & & & \multicolumn{2}{p{0.3\linewidth}}{
        \cellcolor[rgb]{0.9,0.9,0.9}{
            \makecell[{{p{\linewidth}}}]{
                \texttt{\tiny{[GM$|$GM]}}
                \texttt{Annoyance: 8} \\
\texttt{Acceptance: 0} \\
\texttt{Indifference: 1} \\
\texttt{Resentment: 2} \\
            }
        }
    }
    & & \\ \\

    \theutterance \stepcounter{utterance}  
    & & & \multicolumn{2}{p{0.3\linewidth}}{
        \cellcolor[rgb]{0.9,0.9,0.9}{
            \makecell[{{p{\linewidth}}}]{
                \texttt{\tiny{[GM$|$GM]}}
                \texttt{Annoyance: 7} \\
\texttt{Acceptance: 0} \\
\texttt{Indifference: 0} \\
\texttt{Resentment: 3} \\
            }
        }
    }
    & & \\ \\

    \theutterance \stepcounter{utterance}  
    & & & \multicolumn{2}{p{0.3\linewidth}}{
        \cellcolor[rgb]{0.9,0.9,0.9}{
            \makecell[{{p{\linewidth}}}]{
                \texttt{\tiny{[GM$|$GM]}}
                \texttt{game\_result = LOSE} \\
            }
        }
    }
    & & \\ \\

    \theutterance \stepcounter{utterance}  
    & & \multicolumn{4}{p{0.6\linewidth}}{
        \cellcolor[rgb]{0.95,0.95,0.95}{
            \makecell[{{p{\linewidth}}}]{
                \texttt{\tiny{[P1$\leftrightarrow$P]}}
                \texttt{Post{-}episode Reflexion (not sent to GM):} \\
\texttt{New plan: My previous attempts failed because I treated the scoring as a binary choice between "Annoyance" and "Indifference" without fully accounting for the cumulative effect of the doctor's persistence on Olivia's emotional state. While I correctly identified "Annoyance" as the primary emotion in Trial 3, the specific distribution of scores (likely 9 for Annoyance, 1 for Resentment, 0 for others) may have been too rigid or missed the nuance that her "defensive barrier" implies a growing, simmering frustration that borders on resentment, even if it isn't deep{-}seated bitterness. The dialogue shows a progression: she starts with mild dismissal ("I suppose") and escalates to active rejection ("makes things more complicated"). This escalation suggests that while "Annoyance" is the dominant feeling, "Resentment" should be scored higher than 1 (perhaps 3{-}4) to reflect her feeling that the doctor is wasting her time and ignoring her stated boundaries. "Indifference" must remain 0 because she is actively engaging to stop the conversation, and "Acceptance" must be 0. My new strategy is to assign a high score to "Annoyance" (8{-}9) to capture the immediate irritation, a moderate score to "Resentment" (3{-}4) to capture the frustration with the doctor's refusal to listen, and 0 to both "Indifference" and "Acceptance". This accounts for the active, escalating nature of her pushback.} \\
            }
        }
    }
    & \\ \\

\end{supertabular}
}

\end{document}
